古典シャドウ(ランダムな基底での測定)に、古典計算で作った近似状態(prior)を組み合わせて推定の分散を減らす共通ランダム測定(CRM)\cite{vermersch2024}を、フェルミオンのハバード模型に系統的に適用した。単一のパウリ文字列に対する CRM の分散は元論文が厳密に与えており、分散を減らせる割合(利得)は prior の大域的な忠実度ではなく、その観測量に対する prior の外れ $\Delta$ で決まる。本研究は、ハバード模型の現実的な prior について、この外れが観測量ごとにどう振る舞い、なぜそうなるかを調べた。

第一に、系を 256 量子ビットまで広げて prior の大域的な忠実度が崩れても、局所観測量の利得はほとんど変わらないことを示した。第二に、平均場(非制限 Hartree--Fock)prior は、電荷の量(オンサイトの ZZ 相関)では強結合で誤差が $(t/U)^2$ で消えて利得の天井に届く一方、スピンの量では損をすることを示した。その原因を、スピンの向きを選んでしまうことと、隣どうしの量子的な一重項の相関が足りないことの2つに分けた。前者はスピン回転で平均した混合状態 prior で直り、後者は直らない。第三に、スピン射影 HF の効果が系の大きさに反比例して消えることを回転子の描像で定量的に説明した。結合次元を制限した変分 MPS をスピン回転で平均した prior は、電荷とスピンの両方で損をせず、厳密な状態を使わない prior の中で最も良いことを示した。第四に、二重占有のようなパウリ項の和について CRM の分散を厳密に書き下し、利得が項ごとの外れで決まることを示した。さらに、2 次元シリンダー(最大 96 量子ビット)とドープ系での平均場 prior の有効範囲、ショットの配分、損をしない推定量、読み出し誤差に対する脆さと prior 側での補正を調べた。\todo{2次元の再計算の結果が出たら数値を入れて要旨を確定する}
